% Procedencia: results_v3/temporal_validation.csv, temporal_validation_comparacion.csv (code/v3/temporal_validation.py,
% corpus code/v3/temporal_corpus.py); cifras copiadas de results_v3/temporal_validation.md y temporal_validation_tabla.tex.
% Todos los intervalos son bootstrap de familias completas (B=2000, semilla 0), condicionales a las predicciones ajustadas.
% Idioma: ingles, como paper1_manuscript.tex. Tabla: % generated by code/v3/temporal_report.py -- experiment 7 (temporal validation), principal cut T=2010
% requires \usepackage{booktabs,siunitx,multirow}
\begin{table*}[t]
\centering
\caption{Retrospective protocol estimates versus realised forward-in-time performance. SuperCon corpus dated by first report (\num{10215} compositions first reported before 2010, \num{2328} at or after; 76\,\% of the later compositions belong to elemental systems absent before 2010). Estimate = 5-fold out-of-fold predictions using only the pre-2010 data; realised = one model fitted on all pre-2010 data and evaluated on the later compositions. $\Delta$ = estimate $-$ realised; 95\,\% intervals from a bootstrap of whole chemical families (B = 2000), conditional on the fitted predictions. High-$T_c$ threshold = 90th percentile of the training $T_c$ (operator $\geq$). L1 differs from L0 only in the partition (the corpus holds one row per composition). In the \emph{Realised} row the $\Delta$ columns give the 95\,\% interval of the realised value.}
\label{tab:temporal}
\small
\sisetup{retain-explicit-plus}
\begin{tabular}{ll S[table-format=2.2] S[table-format=-2.2] S[table-format=1.3] S[table-format=-1.3] S[table-format=1.2] S[table-format=-1.2] S[table-format=1.2]}
\toprule
Model & Protocol & {MAE (K)} & {$\Delta$MAE} & {$R^2$} & {$\Delta R^2$} & {top-100} & {$\Delta$top-100} & {Spearman} \\
\midrule
\multirow{7}{*}{XGBoost} & \textbf{Realised (train $<$ 2010, test $\geq$ 2010)} & \bfseries 7.43 & {[6.67, 8.27]} & \bfseries 0.328 & {[0.182, 0.469]} & \bfseries 0.17 & {[0.07, 0.29]} & 0.61 \\
 & L0 random $K$-fold & 4.34 & -3.10 & 0.910 & +0.582 & 0.99 & +0.82 & 0.95 \\
 & L1 grouped by composition & 4.18 & -3.25 & 0.912 & +0.584 & 0.97 & +0.80 & 0.95 \\
 & L2 grouped by elemental system & 6.76 & -0.67 & 0.827 & +0.499 & 0.92 & +0.75 & 0.93 \\
 & L3 grouped by host elements & 7.53 & +0.10 & 0.799 & +0.471 & 0.87 & +0.70 & 0.91 \\
 & L4 leave-clusters-out ($k=50$) & 11.35 & +3.92 & 0.604 & +0.276 & 0.54 & +0.37 & 0.83 \\
 & L0 re-weighted to the test mix of new/seen systems & 4.60 & -2.84 & 0.813 & +0.485 & 0.94 & +0.77 & 0.88 \\
\midrule
\multirow{7}{*}{RandomForest} & \textbf{Realised (train $<$ 2010, test $\geq$ 2010)} & \bfseries 7.49 & {[6.72, 8.32]} & \bfseries 0.339 & {[0.205, 0.472]} & \bfseries 0.15 & {[0.06, 0.27]} & 0.61 \\
 & L0 random $K$-fold & 4.79 & -2.70 & 0.898 & +0.559 & 0.99 & +0.84 & 0.95 \\
 & L1 grouped by composition & 4.63 & -2.86 & 0.903 & +0.563 & 0.99 & +0.84 & 0.95 \\
 & L2 grouped by elemental system & 7.06 & -0.43 & 0.817 & +0.478 & 0.85 & +0.70 & 0.92 \\
 & L3 grouped by host elements & 7.96 & +0.46 & 0.780 & +0.441 & 0.69 & +0.54 & 0.91 \\
 & L4 leave-clusters-out ($k=50$) & 11.28 & +3.79 & 0.626 & +0.287 & 0.73 & +0.58 & 0.83 \\
 & L0 re-weighted to the test mix of new/seen systems & 4.95 & -2.54 & 0.794 & +0.454 & 1.00 & +0.85 & 0.87 \\
\bottomrule
\end{tabular}
\end{table*}
 (booktabs, siunitx, multirow).

\subsection{Which retrospective protocol anticipates performance on later superconductors?}
\label{sec:temporal}

Cross-validation protocols are proxies for a question that can be asked directly of SuperCon: how well does a model trained on
everything known up to a year $T$ predict the superconductors first reported at or after $T$? We dated every SuperCon entry by the
year in its bibliographic field (99.7\,\% of rows; the year of the \emph{report}, not of discovery), applied the chemical exclusions of
Sec.~\ref{sec:audit} to the exact formulas, and collapsed the 20\,771 dated $T_c>0$ reports to 12\,543 normalised compositions
(3\,267 elemental systems), each dated by its first report and labelled with the median $T_c$ of its reports. For each cut
$T\in\{1995, 2005, 2010, 2015\}$ we fitted the Dataset-B models on all compositions first reported before $T$ and evaluated them,
frozen, on those reported at or after $T$; independently, using only the pre-$T$ data, we computed the five-fold out-of-fold estimate
that each protocol L0--L4 would have produced at the time. Between 72\,\% and 84\,\% of the later compositions belong to elemental
systems that did not exist in the database before the cut, so the forward-in-time test is dominated by the regime that grouped
validation is meant to emulate.

The realised performance is far below the random-CV figure at every cut (Table~\ref{tab:temporal}, Fig.~\ref{fig:temporal}). At the
principal cut $T=2010$ (10\,215 training compositions, 2\,328 later ones, 76\,\% in new systems) XGBoost reaches MAE
$7.43$\,K [6.67, 8.27] and $R^2=0.33$ [0.18, 0.47]; on the later compositions from \emph{new} systems $R^2$ is $0.19$ [0.01, 0.34]
(MAE $8.47$\,K), on those from systems already seen it is $0.83$ [0.71, 0.89] (MAE $4.07$\,K), and on later re-measurements of
compositions already in the training set the model recovers MAE $2.39$\,K and $R^2=0.80$. Random five-fold CV on the same pre-2010
data estimates MAE $4.34$\,K and $R^2=0.91$: it understates the realised error by $3.10$\,K [2.25, 4.03] and overstates $R^2$ by
$0.58$ [0.44, 0.73]. The $R^2$ overstatement of random CV is significant at all four cuts (+0.28, +0.60, +0.58, +0.42 for 1995,
2005, 2010, 2015); grouping by composition (L1) changes nothing, because the corpus holds one row per composition. Grouping by
elemental system (L2) or by host elements (L3) moves the MAE estimate into the realised interval at 2005 and 2010
($\Delta$MAE $=-0.67$\,K [$-1.83$, 0.43] and $+0.10$\,K [$-1.00$, 1.12] at 2010) but reduces the $R^2$ overstatement only marginally
(+0.50 and +0.47 at 2010). Leave-clusters-out (L4, $k=50$) is the only protocol whose $R^2$ estimate matches the realised value at
1995 ($\Delta R^2=-0.05$ [$-0.12$, 0.02]) and 2015 (+0.11 [$-0.16$, 0.29]), but it is pessimistic in MAE after 2005
(+3.9\,K [2.6, 5.2] at 2010) and still overstates $R^2$ by $0.28$--$0.31$ at 2005--2010. Averaged over the four cuts, the mean
absolute discrepancy in $R^2$ falls monotonically with the strength of the grouping (L0 0.47, L2 0.39, L3 0.35, L4 0.19), while in
MAE the host-element protocol is closest (L3 2.2\,K; L0 3.8\,K; L4 3.5\,K). Re-weighting the random-CV out-of-fold rows to the
new/seen-system mix of the later compositions does not recover the realised value either ($R^2$ 0.81 versus 0.33 at 2010): the
out-of-fold rows from new systems inside a random split are not representative of the systems that are actually discovered later.

Two further observations bound the comparison. First, the screening metric cannot be transferred across time as a precision: the
prevalence of compositions above the training 90th percentile among the later compositions collapses from 9.2\,\% at 1995 to
0.9--1.8\,\% at 2005--2015, so the realised top-100 precision falls from 0.90 to 0.08--0.47 while every protocol estimates
0.36--0.99; as a lift over prevalence the forward-in-time ranking is in fact strong (18--26$\times$ at 2005--2010), and the Spearman
rank correlation degrades more gently (0.95 retrospective versus 0.57--0.78 realised). Second, the realised $R^2$ is fragile:
removing the 526 training compositions whose year came from a fallback dating rule (5\,\% of the training set, mostly 1960s--1970s
alloys and Chevrel phases) shifts the XGBoost predictions for the then-unseen Bi--S--O family from about 15\,K to about 60\,K and drives the realised
$R^2$ at 2010 from $0.33$ to $-0.20$ (MAE 7.4 to 9.1\,K), whereas the random-CV and leave-clusters-out estimates move by less than 0.3\,K in MAE
and 0.02 in $R^2$ (the host-element estimate by 1.1\,K and 0.05). Using the maximum instead of the median $T_c$ per composition, or re-seeding the random and cluster partitions,
leaves all conclusions unchanged (Supplementary Table).

\emph{Answer to the question.} No retrospective protocol reproduces the forward-in-time performance. Random cross-validation is the
worst proxy for $R^2$ at every cut and for MAE at 1995 and 2005; at 2010 the worst MAE proxy is leave-clusters-out
($+3.9$\,K against $-3.1$\,K for random CV), at 2015 random CV is the closest in MAE because the later compositions are
low-$T_c$ and low-variance, and at 1995 random CV is also the closest in top-100 precision (+0.02), where all protocols agree
within their intervals. The
grouped protocols are consistently closer, with the host-element grouping (L3) the best predictor of the realised MAE and
leave-clusters-out (L4) the best predictor of the realised $R^2$, and the direction of their residual error is stable (random CV
optimistic, LOCO pessimistic in MAE after 2005). The realised $R^2$ on new elemental systems ($0.1$--$0.5$) is the number a user of
these models should expect, and none of the protocols, applied to the data available at the time, would have predicted it to
within $0.1$ at more than one cut.

% Figura: results_v3/temporal_validation_fig.pdf
% \begin{figure*}[t]\centering\includegraphics[width=\textwidth]{results_v3/temporal_validation_fig.pdf}
% \caption{Retrospective out-of-fold estimate of each protocol (points, 95\,\% family-bootstrap intervals) against the realised
% forward-in-time performance (line and band) for four cuts; XGBoost in colour, RandomForest in grey (realised value at the right
% of each panel). Open diamonds: random CV re-weighted to the new/seen-system mix of the later compositions.}\label{fig:temporal}
% \end{figure*}
